\begin{mistake}{2026-04-12}{31}

\begin{problemcontent}
设 \( y(x) \) 是微分方程 \( y^{(4)}+y''=0 \) 的不恒等于零的解，且 \( \lim_{x\to 0}y(x)=0 \)。求 \( y(x) \) 关于无穷小量 \( x \) 的最高阶数。
\end{problemcontent}

\begin{solutioncontent}
先求方程的通解。

特征方程为
\[
r^4+r^2=0
\]
即
\[
r^2(r^2+1)=0
\]
故特征根为 \(r=0\)（二重根），\(r=\pm i\)。于是通解为
\[
y=C_1+C_2x+C_3\cos x+C_4\sin x
\]

由 \(\lim_{x\to 0}y(x)=0\) 可知 \(y(0)=0\)，代入得
\[
C_1+C_3=0
\]
即 \(C_1=-C_3\)。因此可写成
\[
y=C_2x+C_4\sin x+C_3(\cos x-1)
\]

在 \(x\to 0\) 时展开：
\[
\sin x=x-\frac{x^3}{6}+o(x^3),\qquad \cos x-1=-\frac{x^2}{2}+\frac{x^4}{24}+o(x^4)
\]
故
\[
y=(C_2+C_4)x-\frac{C_3}{2}x^2-\frac{C_4}{6}x^3+o(x^3)
\]

为了使无穷小阶数尽可能高，应尽量消去低阶项。

1. 令一次项系数为 0：
\[
C_2+C_4=0
\]

2. 再令二次项系数为 0：
\[
C_3=0
\]

此时
\[
y=-\frac{C_4}{6}x^3+o(x^3)
\]
只要取 \(C_4\neq 0\)，就得到非零解，且其阶数为 3。

若还想使阶数高于 3，则还需令三次项系数为 0，即 \(C_4=0\)。这时又由 \(C_2=-C_4=0\)、\(C_3=0\)、\(C_1=-C_3=0\)，只能得到零解，与题设“不恒等于零”矛盾。

因此，\(y(x)\) 关于无穷小量 \(x\) 的最高阶数为
\[
3
\]
\end{solutioncontent}

\mistakeknowledge{
\begin{itemize}
  \item \textbf{核心公式}：由特征方程 \(r^4+r^2=0\) 写出通解 \(y=C_1+C_2x+C_3\cos x+C_4\sin x\)；无穷小阶数由展开式中最低次非零项决定。
  \item \textbf{易错点}：“最高阶数”是问最低非零次幂最多能提高到几次，不是看展开里出现的最高次幂；处理 \(\sin x,\cos x\) 组合时必须保留足够阶的泰勒项，不能只做等价替换。
  \item \textbf{关联知识}：线性微分方程通解含待定常数时，可通过调节常数消去低阶项，从而提高零点阶数，这是零点重数与泰勒展开结合的典型考法。
\end{itemize}}

\end{mistake}
